\documentclass[11pt,a4paper]{article}

\usepackage[utf8]{inputenc}
\usepackage[T1]{fontenc}
\usepackage{amsmath,amssymb,amsthm}
\usepackage{mathtools}
\usepackage{booktabs}
\usepackage{graphicx}
\usepackage{hyperref}
\usepackage[margin=2.5cm]{geometry}
\usepackage{enumitem}
% \usepackage{microtype}
\usepackage{float}

% Theorem environments
\newtheorem{theorem}{Theorem}
\newtheorem{proposition}[theorem]{Proposition}
\newtheorem{lemma}[theorem]{Lemma}
\newtheorem{corollary}[theorem]{Corollary}
\newtheorem{conjecture}[theorem]{Conjecture}
\theoremstyle{definition}
\newtheorem{definition}[theorem]{Definition}
\newtheorem{remark}[theorem]{Remark}

% Notation
\newcommand{\Phi}{\varPhi}
\DeclareMathOperator{\lcm}{lcm}

\title{\textbf{The Cyclotomic Factor Cascade as a\\Supercritical Branching Process}}
\author{Raphael Christi\thanks{Independent researcher, Saquarema, RJ, Brazil. Email: \texttt{bilubiteco@protonmail.com}}}
\date{March 2026}

\begin{document}
\maketitle

\begin{abstract}
We model the factor chain induced by the third cyclotomic polynomial $\Phi_3(p) = p^2 + p + 1$ as a Galton--Watson branching process on the set of odd primes. For each prime~$p$ in a cascade seeded by a finite set~$S_0$, we count the number of \emph{new} odd prime factors of $\Phi_3(p)$ not already present in the running set~$S$. Over $300$ seed configurations generating $12{,}927$ parent--child observations, we find a global mean offspring rate
\[
\bar\mu \;=\; 1.1747,
\]
placing the process firmly in the \textbf{supercritical} regime ($\mu > 1$). We observe a phase transition: $\mu(p) < 1$ for primes $p \lesssim 200$ (subcritical) and $\mu(p) > 1$ for $p \gtrsim 200$ (supercritical), consistent with the prediction of Stewart's theorem that $P^+(\Phi_3(p)) > p^{2-\varepsilon}$ for large~$p$. The offspring rate increases further for higher cyclotomic polynomials ($\mu \approx 2.99$ for $\Phi_5$, $\mu \approx 1.71$ for $\Phi_7$). These measurements provide a quantitative explanation for the empirical observation that no finite set of odd primes is closed under the $\sigma$-induced factor map, offering probabilistic evidence against the existence of odd perfect numbers.
\end{abstract}

\noindent\textbf{Keywords:} odd perfect numbers, cyclotomic polynomials, branching processes, factor chains, Galton--Watson process

\noindent\textbf{MSC 2020:} 11A25 (primary), 11Y70, 60J80 (secondary)

\bigskip

%=====================================================================
\section{Introduction}
%=====================================================================

A natural number $N$ is \emph{perfect} if $\sigma(N) = 2N$, where $\sigma$ denotes the sum-of-divisors function. By the Euclid--Euler theorem, every even perfect number has the form $2^{p-1}(2^p - 1)$ with $2^p - 1$ prime. Whether any odd perfect number (OPN) exists has been open since antiquity.

Euler proved that any OPN must have the form
\begin{equation}\label{eq:euler}
N = q^\alpha \cdot p_1^{2e_1} \cdots p_k^{2e_k},
\end{equation}
where $q \equiv \alpha \equiv 1 \pmod{4}$ is the \emph{special prime} and all $e_i \geq 1$. Modern results impose increasingly severe necessary conditions: $N > 10^{2200}$ \cite{OchemRao2012}, $k \geq 10$ \cite{Nielsen2015}, and largest prime factor exceeding $10^8$ \cite{GotoOhno2006}.

A key technique in these results is the \emph{factor chain} (or \emph{network method}): since $\sigma(N) = 2N$ and $\sigma$ is multiplicative, every odd prime dividing $\sigma(p_i^{2e_i})$ must itself divide~$N$, hence belong to $S = \{q, p_1, \ldots, p_k\}$. This \emph{closure condition} propagates: new primes force new factors of~$\sigma$, generating a cascade. Nielsen \cite{Nielsen2007} and Ochem--Rao \cite{OchemRao2012} exploit this cascade computationally within branching search trees to derive lower bounds.

In the simplest case $e_i = 1$, the relevant factor is $\sigma(p^2) = \Phi_3(p) = p^2 + p + 1$, where $\Phi_3$ is the third cyclotomic polynomial. The cascade then takes the form:
\[
S_{n+1} = S_n \cup \bigl\{\text{odd prime factors of } \Phi_3(p) : p \in S_n \bigr\}.
\]
If this process stabilizes at a finite set, closure is achieved. If it diverges, no finite closed set exists.

\medskip

\noindent\textbf{Our contribution.} We reinterpret this cascade as a \emph{Galton--Watson branching process} and compute the mean offspring rate $\mu(p)$ empirically. To our knowledge, this probabilistic framing---specifically, the computation of $\mu$ as a function of prime size and the identification of a subcritical-to-supercritical phase transition---has not appeared in the literature. While the divergence of factor chains is well known empirically, the branching process framework provides a quantitative explanation grounded in classical probability theory: a Galton--Watson process with $\mu > 1$ survives (diverges) with positive probability, and the extinction probability decreases exponentially with population size \cite{AthreyaNey1972}.

%=====================================================================
\section{Setup and Definitions}
%=====================================================================

\begin{definition}[Cyclotomic cascade]
Let $S_0$ be a finite set of odd primes. Define the \emph{$\Phi_3$-cascade} by
\[
S_{n+1} = S_n \cup \bigl\{q \text{ odd prime} : q \mid \Phi_3(p) \text{ for some } p \in S_n\bigr\}.
\]
Write $S_\infty = \bigcup_{n \geq 0} S_n$. The cascade \emph{closes} if $S_\infty$ is finite.
\end{definition}

\begin{definition}[Offspring count]
For a prime $p$ encountered at round $n$ of the cascade, define its \emph{offspring} as the set of odd primes dividing $\Phi_3(p)$ that are \textbf{not} already in $S_n$ at the time $p$ is processed:
\[
\mathcal{C}(p, S_n) = \bigl\{q \text{ odd prime} : q \mid \Phi_3(p),\; q \notin S_n \bigr\}.
\]
The offspring count is $X_p = |\mathcal{C}(p, S_n)|$.
\end{definition}

In a classical Galton--Watson process, each individual produces offspring independently with a common distribution. Our setting departs from this idealization in two ways: (1)~the offspring distribution depends on the arithmetic properties of $p$ (it is not identically distributed), and (2)~offspring are suppressed if they already appear in~$S_n$ (a depletion effect). Nevertheless, the mean offspring rate $\mu$ retains its classical significance: if $\mu > 1$, the population grows exponentially in expectation.

We also recall a classical fact constraining which primes can appear.

\begin{proposition}[Residue restriction {\cite[Theorem~2]{Nielsen2007}}]\label{prop:residue}
If $q$ is an odd prime dividing $\Phi_n(p)$, then either $q \mid n$ or $q \equiv 1 \pmod{n}$. In particular, odd prime factors of $\Phi_3(p)$ satisfy $q = 3$ or $q \equiv 1 \pmod{3}$.
\end{proposition}

%=====================================================================
\section{Computational Method}
%=====================================================================

We generated three families of seed sets:
\begin{itemize}[nosep]
\item 200 singletons $\{p\}$ for primes $p \equiv 1 \pmod{3}$, $7 \leq p < 10{,}000$;
\item 60 pairs $\{3, p\}$ for such primes $p$;
\item 40 triples $\{3, p_1, p_2\}$ from small primes $\equiv 1 \pmod{3}$.
\end{itemize}
For each of the resulting 300 seeds, we iterated the $\Phi_3$-cascade for up to 12 rounds or until $|S| > 200$. Factorization used trial division up to $10^4$ combined with deterministic Miller--Rabin primality testing (witnesses $\{2, 3, 5, 7, 11, 13, 17, 19, 23, 29, 31, 37\}$, unconditionally correct below $3.3 \times 10^{24}$). For each prime $p$ processed, we recorded its offspring count $X_p$.

Total runtime: 0.5 seconds on an Apple Mac Mini M4, using only the Python standard library. All code is available at \cite{CodeRepo}.

%=====================================================================
\section{Results}
%=====================================================================

\subsection{Global offspring rate}

Over 300 cascades, we recorded $12{,}927$ parent--child observations. The global statistics are:

\begin{table}[H]
\centering
\begin{tabular}{@{}lr@{}}
\toprule
\textbf{Quantity} & \textbf{Value} \\
\midrule
Total parent primes & 12,927 \\
Total offspring generated & 15,187 \\
\textbf{Mean offspring $\bar\mu$} & \textbf{1.1747} \\
Cascades closed ($S_\infty$ finite) & 0 / 300 \\
Cascades diverged & 300 / 300 \\
\bottomrule
\end{tabular}
\caption{Global statistics for the $\Phi_3$ cascade over 300 seed configurations.}
\label{tab:global}
\end{table}

Since $\bar\mu > 1$, the cascade is \textbf{supercritical} in the Galton--Watson sense. No cascade closed.

\subsection{Offspring distribution}

The empirical offspring distribution is concentrated on small values:

\begin{table}[H]
\centering
\begin{tabular}{@{}crc@{}}
\toprule
\textbf{Offspring $k$} & \textbf{Count} & \textbf{Fraction} \\
\midrule
0 & 1{,}351 & 0.105 \\
1 & 8{,}570 & 0.663 \\
2 & 2{,}497 & 0.193 \\
3 & 421 & 0.033 \\
4 & 82 & 0.006 \\
$\geq 5$ & 6 & $< 0.001$ \\
\bottomrule
\end{tabular}
\caption{Offspring distribution. The mode is $k=1$ (66\% of parents produce exactly one new prime).}
\label{tab:offspring}
\end{table}

The distribution resembles a shifted geometric with mean $\mu \approx 1.17$, consistent with a mildly supercritical branching process. The extinction probability for a Galton--Watson process with this offspring distribution, estimated via the generating function $\phi(s) = \sum p_k s^k$, is $\rho \approx 0.73$---high, but strictly less than~$1$.

\subsection{Phase transition: $\mu(p)$ as a function of prime size}\label{sec:phase}

Binning the offspring counts by $\lfloor 2\log_{10} p \rceil / 2$ (bins of width $0.5$ in $\log_{10}$), we observe a clear transition:

\begin{table}[H]
\centering
\begin{tabular}{@{}rrrrr@{}}
\toprule
$\log_{10}(p)$ & $\sim p$ range & $n$ & $\mu(p)$ & Regime \\
\midrule
0.5 & 3--5 & 300 & 0.867 & subcritical \\
1.0 & 6--17 & 592 & 0.840 & subcritical \\
1.5 & 18--56 & 1{,}108 & 0.956 & subcritical \\
2.0 & 56--178 & 2{,}212 & 0.895 & subcritical \\
2.5 & 178--562 & 1{,}880 & 1.138 & \textbf{supercritical} \\
3.0 & 562--1{,}778 & 1{,}802 & 1.407 & \textbf{supercritical} \\
3.5 & 1{,}778--5{,}623 & 1{,}524 & 1.472 & \textbf{supercritical} \\
4.0 & 5{,}623--17{,}783 & 776 & 1.171 & \textbf{supercritical} \\
4.5 & 17{,}783--56{,}234 & 372 & 1.911 & \textbf{supercritical} \\
5.0 & 56{,}234--177{,}828 & 470 & 1.402 & \textbf{supercritical} \\
\bottomrule
\end{tabular}
\caption{Mean offspring $\mu(p)$ by prime size. The transition occurs near $p \approx 200$.}
\label{tab:mu_bins}
\end{table}

The transition from subcritical to supercritical near $\log_{10}(p) \approx 2.3$ (i.e., $p \approx 200$) is consistent with the following theoretical prediction.

\begin{theorem}[Stewart {\cite{Stewart2013}}]\label{thm:stewart}
For $f(x)$ irreducible of degree $d \geq 2$ over $\mathbb{Z}$,
\[
P^+\bigl(f(n)\bigr) > n^{1 + 1/(d-1) - \varepsilon}
\]
for all $n > N_0(\varepsilon)$, where $P^+$ denotes the largest prime factor.
\end{theorem}

For $\Phi_3$ ($d = 2$), this gives $P^+(\Phi_3(p)) > p^{2-\varepsilon}$, meaning the number of digits of the largest prime factor roughly doubles each generation. This forces $\mu(p) \geq 1$ for large $p$, since $\Phi_3(p) \approx p^2$ must have at least one prime factor of size $\gg p$, which is necessarily new.

\subsection{Residue class dependence}

Consistent with Proposition~\ref{prop:residue}:

\begin{table}[H]
\centering
\begin{tabular}{@{}lcr@{}}
\toprule
$p \bmod 3$ & $\mu$ & $n$ \\
\midrule
$\equiv 0$ & 0.867 & 300 \\
$\equiv 1$ & 1.182 & 12{,}627 \\
\bottomrule
\end{tabular}
\caption{Offspring rate by residue class. Primes $\equiv 2 \pmod{3}$ do not appear as factors of $\Phi_3$.}
\label{tab:residue}
\end{table}

The value $\mu \approx 0.87$ for $p = 3$ reflects the special role of $3$ (which divides $\Phi_3(p)$ whenever $p \equiv 1 \pmod{3}$, providing no new information). The bulk of the cascade consists of primes $\equiv 1 \pmod{3}$, for which $\mu = 1.182 > 1$.

\subsection{Largest child ratio and connection to Stewart}

For each prime $p$ in the cascade, we computed the ratio $\log_{10}(\text{largest child}) / \log_{10}(p)$. Over 300 primes with $p > 10$:
\[
\text{mean ratio} = 1.339, \qquad \text{median ratio} = 1.345.
\]
Stewart's theorem predicts this ratio tends to $2.0$ for large $p$. The observed value $\approx 1.34$ for primes up to $\sim 10^4$ suggests the Stewart regime is not yet fully reached, but the trend is upward, consistent with $\mu(p)$ increasing with $p$.

\subsection{Higher cyclotomic polynomials}

For a non-special prime $p$ with exponent $2e$ in an OPN, the relevant factor is $\sigma(p^{2e}) = \prod_{d \mid (2e+1)} \Phi_d(p)$. The dominant new factor is $\Phi_{2e+1}(p)$ of degree $\varphi(2e+1) \geq 2$ for $e \geq 1$. We measured $\mu$ for the first offspring generation of $\Phi_5$ and $\Phi_7$:

\begin{table}[H]
\centering
\begin{tabular}{@{}lccc@{}}
\toprule
& Degree & $\mu$ & Stewart exponent \\
\midrule
$\Phi_3(p) = p^2 + p + 1$ & 2 & 2.987 & $1 + 1/1 = 2.0$ \\
$\Phi_5(p) = p^4 + \cdots + 1$ & 4 & 2.340 & $1 + 1/3 \approx 1.33$ \\
$\Phi_7(p) = p^6 + \cdots + 1$ & 6 & 1.707 & $1 + 1/5 = 1.2$ \\
\bottomrule
\end{tabular}
\caption{Mean offspring for higher cyclotomic polynomials (first generation, 150 primes $\equiv 1 \pmod{3}$).}
\label{tab:multi}
\end{table}

All three are supercritical. Since $\varphi(2e+1) \geq 2$ for all $e \geq 1$, Stewart's theorem applies to every non-special exponent. In a realistic OPN with multiple exponents active simultaneously, the effective $\mu$ is the sum over all active $\Phi_n$, making the cascade diverge even faster.

%=====================================================================
\section{Implications for Odd Perfect Numbers}
%=====================================================================

The closure condition for an OPN requires $S_\infty = S_0$---the cascade must stabilize immediately. In the branching process framework, this corresponds to \emph{immediate extinction}: every individual in the first generation must be sterile ($X_p = 0$ for all $p \in S_0$).

For a set $S_0$ of $k$ primes, the probability of immediate extinction under our empirical offspring distribution is approximately $(p_0)^k = (0.105)^k$, which is $\approx 10^{-10}$ for $k = 10$ (the minimum number of distinct prime factors required by Nielsen~\cite{Nielsen2015}). But immediate extinction is necessary \emph{at every round}, not just the first. The probability of the cascade closing within $r$ rounds is bounded by $\rho^{|S_r|}$, where $\rho \approx 0.73$ is the per-individual extinction probability. Since $|S_r|$ grows exponentially when $\mu > 1$, the survival probability approaches~$1$ exponentially fast.

\begin{remark}
This is a heuristic argument, not a proof. The Galton--Watson model assumes independence of offspring, which is only approximately true in our setting (the residue restriction of Proposition~\ref{prop:residue} introduces correlations). A rigorous proof would require showing that the correlations do not reduce $\mu$ below~$1$---equivalently, that the depletion effect (primes already in $S$) does not dominate the growth. Our data suggest it does not: the depletion effect is strongest for small primes (where $\mu < 1$), but these are outnumbered by large primes (where $\mu > 1$) as the cascade progresses.
\end{remark}

\begin{remark}
Pomerance's heuristic \cite{Pomerance2003} argues against OPNs via the global probability that $\sigma(N) = 2N$ for a random odd number of the appropriate form. Our approach is complementary: we argue against OPNs via the \emph{local} branching rate of the factor cascade. The two perspectives are independent and reinforce each other.
\end{remark}

%=====================================================================
\section{Relation to Prior Work}
%=====================================================================

The factor chain method originates with Sylvester \cite{Sylvester1888}, who manually traced small cascades. Nielsen \cite{Nielsen2007} systematized this into the ``network method'' used to prove $k \geq 9$, later extended to $k \geq 10$ \cite{Nielsen2015}. Ochem and Rao \cite{OchemRao2012} pushed the lower bound on $N$ to $10^{1500}$ (later $10^{2200}$) using branching search trees built on factor chains. Fletcher, Nielsen, and Ochem \cite{FletcherNielsenOchem} combined sieve methods with factor chains.

Our contribution is not the cascade itself, which is well established, but the \emph{probabilistic framing}:
\begin{enumerate}[nosep]
\item Computing the empirical offspring distribution and mean $\mu$ across $12{,}927$ observations;
\item Identifying the subcritical-to-supercritical phase transition near $p \approx 200$;
\item Connecting $\mu(p) > 1$ to Stewart's theorem as the theoretical explanation;
\item Measuring $\mu$ for $\Phi_5$ and $\Phi_7$, showing universality of supercriticality.
\end{enumerate}

The BYU group's work on spoof perfect factorizations \cite{NielsenSpoofs} proved that infinitely many spoofs exist, implying that approaches using only the multiplicative structure of $\sigma$ cannot resolve the OPN problem. Our results are consistent with this: the branching process model captures the multiplicative structure and shows it is generically divergent, but cannot distinguish primes from composites (the essential difference between real OPNs and spoofs).

%=====================================================================
\section{Conclusion}
%=====================================================================

We have shown that the cyclotomic factor cascade, when modeled as a Galton--Watson branching process, is supercritical with mean offspring $\bar\mu = 1.175$. This provides a quantitative explanation for the empirical observation---noted implicitly by every computational study of OPNs since Sylvester---that factor chains diverge. The phase transition at $p \approx 200$, the consistency with Stewart's theorem, and the even higher $\mu$ for $\Phi_5$ and $\Phi_7$ all point to the same conclusion: the closure condition required for an OPN is a probability-zero event in the branching process model.

A natural next step would be to make this rigorous: prove that for any finite set $S$ of odd primes with $|S| \geq k_0$, the expected number of new primes generated by one round of the $\Phi_3$-cascade exceeds $|S|$ (i.e., the process is supercritical in a formal sense). Combined with Stewart's theorem for large primes and computational verification for small primes, this could yield a conditional proof that no finite closed set exists.

\bigskip

\begin{thebibliography}{99}

\bibitem{AthreyaNey1972}
K.~B.~Athreya and P.~E.~Ney,
\emph{Branching Processes},
Springer, 1972.

\bibitem{FletcherNielsenOchem}
S.~A.~Fletcher, P.~P.~Nielsen, and P.~Ochem,
Sieve methods for odd perfect numbers,
\emph{Math.\ Comp.}\ \textbf{81} (2012), no.~279, 1753--1776.

\bibitem{GotoOhno2006}
T.~Goto and Y.~Ohno,
Odd perfect numbers have a prime factor exceeding $10^8$,
\emph{Math.\ Comp.}\ \textbf{77} (2008), no.~263, 1859--1868.

\bibitem{Nielsen2007}
P.~P.~Nielsen,
Odd perfect numbers have at least nine distinct prime factors,
\emph{Math.\ Comp.}\ \textbf{76} (2007), no.~260, 2109--2126.

\bibitem{Nielsen2015}
P.~P.~Nielsen,
Odd perfect numbers, Diophantine equations, and upper bounds,
\emph{Math.\ Comp.}\ \textbf{84} (2015), no.~295, 2549--2567.

\bibitem{NielsenSpoofs}
P.~P.~Nielsen \emph{et al.},
Odd, spoof perfect factorizations,
\emph{J.\ Number Theory}\ \textbf{234} (2022), 31--47.

\bibitem{OchemRao2012}
P.~Ochem and M.~Rao,
Odd perfect numbers are greater than $10^{1500}$,
\emph{Math.\ Comp.}\ \textbf{81} (2012), no.~279, 1869--1877.

\bibitem{Pomerance2003}
C.~Pomerance,
Personal communication to J.~Zelinsky, 2003;
see also \emph{Multiply perfect numbers, Mersenne primes, and effective computability},
\emph{Math.\ Ann.}\ \textbf{226} (1977), no.~3, 195--206.

\bibitem{Stewart2013}
C.~L.~Stewart,
On the greatest prime factor of integers of the form $ab + 1$,
\emph{Acta Arith.}\ \textbf{163} (2014), no.~4, 299--331.

\bibitem{Sylvester1888}
J.~J.~Sylvester,
Sur les nombres parfaits,
\emph{Comptes Rendus}\ \textbf{106} (1888), 403--405.

\bibitem{CodeRepo}
R.~Christi,
Computational scripts for this paper,
available at \url{https://github.com/raphaelchristi/galton-watson-opn}, 2026.

\end{thebibliography}

\end{document}
